\paragraph{Rows and complete compact fields.}
Let $m,W,s$ denote the complex network constants and keep the layer's global
$p,d,K$ at every node. Put
\[
 G=4\lceil\log_2p\rceil+6,\quad H=dG,\quad
 q_F=\left\lceil\frac{2H}{K}\right\rceil,\quad
 q_B=\left\lceil\frac{H}{K}\right\rceil,
\]
\[
 q_0=\lceil\log_2W\rceil\lceil\log_m(2d)\rceil,
 \quad k_0=\lceil\log_m d\rceil,\quad Q_0=q_0+q_F+q_B.
\]
If $D\le Q_0$, apply all $D$ selected kernels individually. Otherwise,
process individually the first $q_0+q_F$ and last $q_B$ selected axes.
This uses existing coordinates only. The remaining active axes are a
consecutive middle block of length $D-Q_0$.

Regard the first $q_0$ chunks as the row index. Immediately after it, the
next $q_F$ chunks have at least $2H$ bits. Reinterpret these consecutive
bits as two $H$-bit front fields and a spectator remainder; a remainder
of length zero is omitted. The last $q_B$ chunks similarly supply one
$H$-bit back field and a spectator remainder. These reinterpretations move
no records and preserve the total address space. For a call with $f$ axes
per slot, carve $nG=(f-1)G\le H$ consecutive bits from each of these
fields; collapse their unused bits into spectators. The resulting field
list has a fixed number of entries. Both front fields precede all active
slots and the back field follows them, so the compact-control proposition
applies to every ordered pair of slots. For $f=1$ no compact fields are used.

The original row range has length $R_{\rm row}=2^{q_0K}\ge W^{k_0}$.
For each original prefix, pad this row range to the next multiple of
$W^{k_0}$ by adding zero rows. A row includes the entire complete front
fields, active middle block, back field, all spectator coordinates and
polynomial record suffix. Padding adds whole such rows, so increases total
volume by at most two and does not change any within-row range.

At a node write the row index as $u=Wg+w$, $0\le w<W$, and send the entire
row to role $w$. The index $u$ consists solely of the first row region:
\emph{no compact-field bit is included in it}. Each role therefore has
the same complete within-row address set and exactly $1/W$ of the parent
volume. At depth $j$ the remaining row count is divisible by $W^{k_0-j}$.
This invariant survives further splits, padding, and reassembly. Pointwise
gates on different roles have aligned $g$ and every remaining coordinate.
All compact operations restore the temporary fields before a scalar gate,
recursive child, or returned result uses that layout.

The fixed-tape scheduler is unchanged: park only the other role streams at
the active stack top, execute the child on a complete role stream, and restore
the parked streams on return. Every child has its own current logical volume;
copying, parking, rewinds and clearing cost a fixed multiple of that volume.
The constant number of compact operations per edge use fixed work areas.
Descriptors include the global $p,d,K,G,H$ and the carved field lengths;
all have $O(p)$ bits and there are a fixed number of fields. The existing
record suffix absorbs polynomial setup at every node and per record.

Each completed invocation acts as the desired $C$ tensor on its active axes
and as identity on the reserved and spectator coordinates. It therefore
commutes with the earlier individual kernels, including those on bits now
interpreted as compact fields. Restoring the row permutation returns the
result to each original row separately. A padded zero row is consequently
zero at completion and can be deleted exactly as before. Neither temporary
field contents nor intermediate scratch payloads were assumed to be zero.

The number of individually processed axes is at most $Q_0$, including the
fallback $D\le Q_0$. For large $p$, $K\ge d^c/2$ and
\[
 q_F+q_B\le\frac{3dG}{K}+2\le6d^{1-c}G+2.
\]
Thus preprocessing, padding and reservation cost
\[
 O\bigl(V(\log(2d)+d^{\max\{1-c,0\}}\log p+1)\bigr).
\]
All reserved axes lie within the original $D$ axes; in particular their
total number never exceeds $d$. This fact also bounds their coefficient
dependency depth by the same $8d$ allowance used in the guard proof.

\paragraph{Unrolling the compact-control recurrence.}
Fix rational $0<\beta<1$. Choose $0<\tau<1$ and
$\max\{0,\log_m(s/W)\}<\sigma<1$. Stop when $e<d^\beta$ and apply
the remaining kernels individually. With $F(e)$ denoting time per current
logical volume, the exact row split gives
\[
 F(e)\le(s/W)F(e/m)+O(G^\tau e^\tau+1),\qquad
 F(e)=O(e)\quad(e<d^\beta).
\]
Here $G$ and the compact reservations are global parameters, not rebuilt
at each node. This is why no factor $K^\tau$ appears in the overhead.
Repair costs are already included in the node overhead.

Write $b=s/W$. At level $j$ the internal contribution is bounded by
$O(G^\tau e^\tau(b/m^\tau)^j)$. Summing through the stopping depth gives
\[
 O(G^\tau\log(2d)d^\chi),\qquad
 \chi=\tau+(1-\beta)\max\{\sigma-\tau,0\}.
\]
Indeed if $\sigma\le\tau$ the root bounds every summand; otherwise the
last internal level costs at most the root bound times a fixed constant
and $d^{(1-\beta)(\sigma-\tau)}$. The additive $1$ is absorbed by
$G^\tau e^\tau$ at internal nodes. The leaf contribution is at most
$O(d^{\sigma+\beta(1-\sigma)})$, by the original volume-weighted argument.

Partition $D-Q_0$ into consecutive powers of $m$, using its base-$m$
expansion. There are $O(\log(2d))$ pieces. Include their costs and the
reservation cost to obtain, per original volume,
\[
 O\!\left((\log p)^3
 [d^\chi+d^{\sigma+\beta(1-\sigma)}+d^{\max\{1-c,0\}}+1]\right).
\]
We used $\log d=\Theta(\log p)$ and $G=O(\log p)$ for each fixed
$d=\Theta(p^\epsilon)$ family. Any fixed rational exponents satisfying
\[
 \max\{\tau,\sigma,\chi\}<\lambda<\lambda'<1,\qquad
 \max\{\sigma+\beta(1-\sigma),1-c,0\}<\lambda'
\]
therefore give the complete $O(Vd^{\lambda'})$ layer bound. The strict gaps
absorb the logarithms. The small-$D$ fallback obeys the same bound.
For fixed rational $\beta=u/v$, implement stopping by the integer test
$e^v<d^u$; its polynomial descriptor cost is already included.
